% ===========================================================================
% Experiment Matrix Flow Diagram
% ===========================================================================
% Included by: % ===========================================================================
% Experiment Matrix Flow Diagram
% ===========================================================================
% Included by: % ===========================================================================
% Experiment Matrix Flow Diagram
% ===========================================================================
% Included by: % ===========================================================================
% Experiment Matrix Flow Diagram
% ===========================================================================
% Included by: \input{figures/matrix-flow}
% Shows: the 22 × 4 × 3 × 10 experimental design
%
\begin{tikzpicture}[
  box/.style={
    rectangle,
    draw=black!70,
    fill=#1,
    rounded corners=2pt,
    minimum height=0.7cm,
    font=\small\sffamily,
    align=center,
    thick,
    inner sep=4pt,
  },
  box/.default=blue!8,
  bigbox/.style={
    rectangle,
    draw=black!50,
    fill=#1,
    rounded corners=4pt,
    minimum height=1.0cm,
    font=\sffamily,
    align=center,
    thick,
    inner sep=6pt,
  },
  bigbox/.default=white,
  arrow/.style={-Stealth, thick, black!70},
  multiply/.style={
    circle,
    draw=black!60,
    fill=gray!10,
    font=\small\bfseries,
    inner sep=2pt,
    minimum size=0.5cm,
  },
  total/.style={
    rectangle,
    draw=black!80,
    fill=red!10,
    rounded corners=3pt,
    minimum height=0.8cm,
    font=\small\sffamily\bfseries,
    align=center,
    thick,
    inner sep=5pt,
  },
  sublabel/.style={
    font=\scriptsize\sffamily,
    text=black!60,
  },
]

% ── Column 1: Scenarios ──
\node[box=blue!10] (scenarios) at (0, 0) {22 Scenarios};
\node[sublabel, below=1pt of scenarios] {6 categories};

% ── Multiply ──
\node[multiply] (x1) at (1.8, 0) {$\times$};

% ── Column 2: Models ──
\node[box=green!10] (models) at (3.4, 0) {4 Models};
\node[sublabel, below=1pt of models] {3 providers};

% ── Multiply ──
\node[multiply] (x2) at (5.0, 0) {$\times$};

% ── Column 3: Conditions ──
\node[box=orange!10] (conditions) at (6.8, 0) {3 Conditions};
\node[sublabel, below=1pt of conditions] {B / H / A};

% ── Multiply ──
\node[multiply] (x3) at (8.5, 0) {$\times$};

% ── Column 4: Repetitions ──
\node[box=purple!10] (reps) at (10.2, 0) {10 Reps};
\node[sublabel, below=1pt of reps] {statistical power};

% ── Arrows ──
\draw[arrow] (scenarios) -- (x1);
\draw[arrow] (x1) -- (models);
\draw[arrow] (models) -- (x2);
\draw[arrow] (x2) -- (conditions);
\draw[arrow] (conditions) -- (x3);
\draw[arrow] (x3) -- (reps);

% ── Sub-matrices (below) ──
\node[bigbox=blue!5] (main) at (1.7, -2.2) {
  \textbf{Main Matrix}\\
  {\small $22 \times 4 \times 2 \times 10$}\\
  {\small = 1,760 cells}
};

\node[bigbox=orange!5] (ablation) at (5.2, -2.2) {
  \textbf{Ablation Matrix}\\
  {\small $22 \times 4 \times 7 \times 10$}\\
  {\small = 6,160 cells}
};

\node[bigbox=green!5] (compound) at (9.0, -2.2) {
  \textbf{Compound Matrix}\\
  {\small $15 \times 4 \times 2 \times 10$}\\
  {\small = 1,200 cells}
};

% ── Convergence arrow ──
\node[total] (total) at (5.2, -4.2) {Total: 9,120 experiment cells};

\draw[arrow] (main.south) -- (total.north west);
\draw[arrow] (ablation.south) -- (total.north);
\draw[arrow] (compound.south) -- (total.north east);

% ── Top connection to sub-matrices ──
\draw[arrow, dashed, black!40] (reps.south) -- ++(0, -0.5) -| (main.north);
\draw[arrow, dashed, black!40] (reps.south) -- ++(0, -0.5) -| (ablation.north);
\draw[arrow, dashed, black!40] (reps.south) -- ++(0, -0.5) -| (compound.north);

\end{tikzpicture}

% Shows: the 22 × 4 × 3 × 10 experimental design
%
\begin{tikzpicture}[
  box/.style={
    rectangle,
    draw=black!70,
    fill=#1,
    rounded corners=2pt,
    minimum height=0.7cm,
    font=\small\sffamily,
    align=center,
    thick,
    inner sep=4pt,
  },
  box/.default=blue!8,
  bigbox/.style={
    rectangle,
    draw=black!50,
    fill=#1,
    rounded corners=4pt,
    minimum height=1.0cm,
    font=\sffamily,
    align=center,
    thick,
    inner sep=6pt,
  },
  bigbox/.default=white,
  arrow/.style={-Stealth, thick, black!70},
  multiply/.style={
    circle,
    draw=black!60,
    fill=gray!10,
    font=\small\bfseries,
    inner sep=2pt,
    minimum size=0.5cm,
  },
  total/.style={
    rectangle,
    draw=black!80,
    fill=red!10,
    rounded corners=3pt,
    minimum height=0.8cm,
    font=\small\sffamily\bfseries,
    align=center,
    thick,
    inner sep=5pt,
  },
  sublabel/.style={
    font=\scriptsize\sffamily,
    text=black!60,
  },
]

% ── Column 1: Scenarios ──
\node[box=blue!10] (scenarios) at (0, 0) {22 Scenarios};
\node[sublabel, below=1pt of scenarios] {6 categories};

% ── Multiply ──
\node[multiply] (x1) at (1.8, 0) {$\times$};

% ── Column 2: Models ──
\node[box=green!10] (models) at (3.4, 0) {4 Models};
\node[sublabel, below=1pt of models] {3 providers};

% ── Multiply ──
\node[multiply] (x2) at (5.0, 0) {$\times$};

% ── Column 3: Conditions ──
\node[box=orange!10] (conditions) at (6.8, 0) {3 Conditions};
\node[sublabel, below=1pt of conditions] {B / H / A};

% ── Multiply ──
\node[multiply] (x3) at (8.5, 0) {$\times$};

% ── Column 4: Repetitions ──
\node[box=purple!10] (reps) at (10.2, 0) {10 Reps};
\node[sublabel, below=1pt of reps] {statistical power};

% ── Arrows ──
\draw[arrow] (scenarios) -- (x1);
\draw[arrow] (x1) -- (models);
\draw[arrow] (models) -- (x2);
\draw[arrow] (x2) -- (conditions);
\draw[arrow] (conditions) -- (x3);
\draw[arrow] (x3) -- (reps);

% ── Sub-matrices (below) ──
\node[bigbox=blue!5] (main) at (1.7, -2.2) {
  \textbf{Main Matrix}\\
  {\small $22 \times 4 \times 2 \times 10$}\\
  {\small = 1,760 cells}
};

\node[bigbox=orange!5] (ablation) at (5.2, -2.2) {
  \textbf{Ablation Matrix}\\
  {\small $22 \times 4 \times 7 \times 10$}\\
  {\small = 6,160 cells}
};

\node[bigbox=green!5] (compound) at (9.0, -2.2) {
  \textbf{Compound Matrix}\\
  {\small $15 \times 4 \times 2 \times 10$}\\
  {\small = 1,200 cells}
};

% ── Convergence arrow ──
\node[total] (total) at (5.2, -4.2) {Total: 9,120 experiment cells};

\draw[arrow] (main.south) -- (total.north west);
\draw[arrow] (ablation.south) -- (total.north);
\draw[arrow] (compound.south) -- (total.north east);

% ── Top connection to sub-matrices ──
\draw[arrow, dashed, black!40] (reps.south) -- ++(0, -0.5) -| (main.north);
\draw[arrow, dashed, black!40] (reps.south) -- ++(0, -0.5) -| (ablation.north);
\draw[arrow, dashed, black!40] (reps.south) -- ++(0, -0.5) -| (compound.north);

\end{tikzpicture}

% Shows: the 22 × 4 × 3 × 10 experimental design
%
\begin{tikzpicture}[
  box/.style={
    rectangle,
    draw=black!70,
    fill=#1,
    rounded corners=2pt,
    minimum height=0.7cm,
    font=\small\sffamily,
    align=center,
    thick,
    inner sep=4pt,
  },
  box/.default=blue!8,
  bigbox/.style={
    rectangle,
    draw=black!50,
    fill=#1,
    rounded corners=4pt,
    minimum height=1.0cm,
    font=\sffamily,
    align=center,
    thick,
    inner sep=6pt,
  },
  bigbox/.default=white,
  arrow/.style={-Stealth, thick, black!70},
  multiply/.style={
    circle,
    draw=black!60,
    fill=gray!10,
    font=\small\bfseries,
    inner sep=2pt,
    minimum size=0.5cm,
  },
  total/.style={
    rectangle,
    draw=black!80,
    fill=red!10,
    rounded corners=3pt,
    minimum height=0.8cm,
    font=\small\sffamily\bfseries,
    align=center,
    thick,
    inner sep=5pt,
  },
  sublabel/.style={
    font=\scriptsize\sffamily,
    text=black!60,
  },
]

% ── Column 1: Scenarios ──
\node[box=blue!10] (scenarios) at (0, 0) {22 Scenarios};
\node[sublabel, below=1pt of scenarios] {6 categories};

% ── Multiply ──
\node[multiply] (x1) at (1.8, 0) {$\times$};

% ── Column 2: Models ──
\node[box=green!10] (models) at (3.4, 0) {4 Models};
\node[sublabel, below=1pt of models] {3 providers};

% ── Multiply ──
\node[multiply] (x2) at (5.0, 0) {$\times$};

% ── Column 3: Conditions ──
\node[box=orange!10] (conditions) at (6.8, 0) {3 Conditions};
\node[sublabel, below=1pt of conditions] {B / H / A};

% ── Multiply ──
\node[multiply] (x3) at (8.5, 0) {$\times$};

% ── Column 4: Repetitions ──
\node[box=purple!10] (reps) at (10.2, 0) {10 Reps};
\node[sublabel, below=1pt of reps] {statistical power};

% ── Arrows ──
\draw[arrow] (scenarios) -- (x1);
\draw[arrow] (x1) -- (models);
\draw[arrow] (models) -- (x2);
\draw[arrow] (x2) -- (conditions);
\draw[arrow] (conditions) -- (x3);
\draw[arrow] (x3) -- (reps);

% ── Sub-matrices (below) ──
\node[bigbox=blue!5] (main) at (1.7, -2.2) {
  \textbf{Main Matrix}\\
  {\small $22 \times 4 \times 2 \times 10$}\\
  {\small = 1,760 cells}
};

\node[bigbox=orange!5] (ablation) at (5.2, -2.2) {
  \textbf{Ablation Matrix}\\
  {\small $22 \times 4 \times 7 \times 10$}\\
  {\small = 6,160 cells}
};

\node[bigbox=green!5] (compound) at (9.0, -2.2) {
  \textbf{Compound Matrix}\\
  {\small $15 \times 4 \times 2 \times 10$}\\
  {\small = 1,200 cells}
};

% ── Convergence arrow ──
\node[total] (total) at (5.2, -4.2) {Total: 9,120 experiment cells};

\draw[arrow] (main.south) -- (total.north west);
\draw[arrow] (ablation.south) -- (total.north);
\draw[arrow] (compound.south) -- (total.north east);

% ── Top connection to sub-matrices ──
\draw[arrow, dashed, black!40] (reps.south) -- ++(0, -0.5) -| (main.north);
\draw[arrow, dashed, black!40] (reps.south) -- ++(0, -0.5) -| (ablation.north);
\draw[arrow, dashed, black!40] (reps.south) -- ++(0, -0.5) -| (compound.north);

\end{tikzpicture}

% Shows: the 22 × 4 × 3 × 10 experimental design
%
\begin{tikzpicture}[
  box/.style={
    rectangle,
    draw=black!70,
    fill=#1,
    rounded corners=2pt,
    minimum height=0.7cm,
    font=\small\sffamily,
    align=center,
    thick,
    inner sep=4pt,
  },
  box/.default=blue!8,
  bigbox/.style={
    rectangle,
    draw=black!50,
    fill=#1,
    rounded corners=4pt,
    minimum height=1.0cm,
    font=\sffamily,
    align=center,
    thick,
    inner sep=6pt,
  },
  bigbox/.default=white,
  arrow/.style={-Stealth, thick, black!70},
  multiply/.style={
    circle,
    draw=black!60,
    fill=gray!10,
    font=\small\bfseries,
    inner sep=2pt,
    minimum size=0.5cm,
  },
  total/.style={
    rectangle,
    draw=black!80,
    fill=red!10,
    rounded corners=3pt,
    minimum height=0.8cm,
    font=\small\sffamily\bfseries,
    align=center,
    thick,
    inner sep=5pt,
  },
  sublabel/.style={
    font=\scriptsize\sffamily,
    text=black!60,
  },
]

% ── Column 1: Scenarios ──
\node[box=blue!10] (scenarios) at (0, 0) {22 Scenarios};
\node[sublabel, below=1pt of scenarios] {6 categories};

% ── Multiply ──
\node[multiply] (x1) at (1.8, 0) {$\times$};

% ── Column 2: Models ──
\node[box=green!10] (models) at (3.4, 0) {4 Models};
\node[sublabel, below=1pt of models] {3 providers};

% ── Multiply ──
\node[multiply] (x2) at (5.0, 0) {$\times$};

% ── Column 3: Conditions ──
\node[box=orange!10] (conditions) at (6.8, 0) {3 Conditions};
\node[sublabel, below=1pt of conditions] {B / H / A};

% ── Multiply ──
\node[multiply] (x3) at (8.5, 0) {$\times$};

% ── Column 4: Repetitions ──
\node[box=purple!10] (reps) at (10.2, 0) {10 Reps};
\node[sublabel, below=1pt of reps] {statistical power};

% ── Arrows ──
\draw[arrow] (scenarios) -- (x1);
\draw[arrow] (x1) -- (models);
\draw[arrow] (models) -- (x2);
\draw[arrow] (x2) -- (conditions);
\draw[arrow] (conditions) -- (x3);
\draw[arrow] (x3) -- (reps);

% ── Sub-matrices (below) ──
\node[bigbox=blue!5] (main) at (1.7, -2.2) {
  \textbf{Main Matrix}\\
  {\small $22 \times 4 \times 2 \times 10$}\\
  {\small = 1,760 cells}
};

\node[bigbox=orange!5] (ablation) at (5.2, -2.2) {
  \textbf{Ablation Matrix}\\
  {\small $22 \times 4 \times 7 \times 10$}\\
  {\small = 6,160 cells}
};

\node[bigbox=green!5] (compound) at (9.0, -2.2) {
  \textbf{Compound Matrix}\\
  {\small $15 \times 4 \times 2 \times 10$}\\
  {\small = 1,200 cells}
};

% ── Convergence arrow ──
\node[total] (total) at (5.2, -4.2) {Total: 9,120 experiment cells};

\draw[arrow] (main.south) -- (total.north west);
\draw[arrow] (ablation.south) -- (total.north);
\draw[arrow] (compound.south) -- (total.north east);

% ── Top connection to sub-matrices ──
\draw[arrow, dashed, black!40] (reps.south) -- ++(0, -0.5) -| (main.north);
\draw[arrow, dashed, black!40] (reps.south) -- ++(0, -0.5) -| (ablation.north);
\draw[arrow, dashed, black!40] (reps.south) -- ++(0, -0.5) -| (compound.north);

\end{tikzpicture}
